% Book 3, Year 1, chapter 25 -- shortest paths.
% Provenance (never printed): sources/book3/BRIEFS.md, chapter 25.
% ledger: b1:dijkstra-1959, b1:bellman-1958, b1:floyd-1962, b1:warshall-1962
\chapter{Shortest Paths}\label{ch:b1:shortest-paths}

From the Town Square to the Hospital, the shortest way is 6.3 kilometres long and takes nine minutes at the speed
limits. Another way is a kilometre longer and takes a minute and a half less, because it uses the faster roads.
\omterm{def:b1:graph-traversals:bfs}{Breadth-first search} counts roads, neither metres nor seconds. This chapter weighs every road with a number, here its
travel time or its length, and finds the paths of least total weight: from one place with Dijkstra's algorithm when
weights are not negative, with Bellman--Ford's when they may be, in linear time when the \omterm{def:b1:graph-traversals:graph}{graph} is acyclic, and between
all pairs of places at once with Floyd--Warshall's.

\begin{recall}
The second book ran Dijkstra's algorithm by hand on small maps. \omterm{def:b1:heaps:binary-heap}{Binary heaps} and \omterm{def:b1:heaps:pq}{priority queues} were built in
\cref{ch:b1:heaps}; \omterm{def:b1:graph-traversals:topo}{topological orders} and \omterm{def:b1:graph-traversals:bfs}{breadth-first search} in \cref{ch:b1:graph-traversals}; tables of
subproblems in \cref{ch:b1:dynamic-programming}.
\end{recall}

\section{Weighted shortest paths}\label{sec:b1:shortest-paths:weighted}

\begin{definition}[Weight of a path, shortest path, weighted distance]\label{def:b1:shortest-paths:shortest}
In a \omterm{def:b1:graph-traversals:graph}{directed graph} with a weight $w(u, v)$ on each \omterm{def:b1:graph-traversals:graph}{arc}, the \emph{weight} of a path is the sum of the weights of its
\omterm{def:b1:graph-traversals:graph}{arcs}. The \emph{weighted distance}\index{weighted distance} $\delta(s, v)$ is the infimum of the weights of the paths
from $s$ to $v$: $+\infty$ if there is none, $-\infty$ if there are paths of arbitrarily small weight. A
\emph{shortest path}\index{shortest path} from $s$ to $v$ is a path of weight $\delta(s, v)$.
\end{definition}

An undirected road is two \omterm{def:b1:graph-traversals:graph}{arcs} of the same weight. The town's roads carry their length in metres, or their travel time
at the speed limit, rounded to the second.

\omcode{code/bachelor-1/25-shortest-paths/listings/wgraph.py}{14}{25}{The roads weighted by travel
time}\label{lst:b1:shortest-paths:wgraph}

\omoutput{code/bachelor-1/25-shortest-paths/outputs/wgraph.txt}{The road from the Station to the Town Hall is short,
1000 m, but limited to 30 km/h.}\label{out:b1:shortest-paths:wgraph}

\begin{definition}[Negative cycle]\label{def:b1:shortest-paths:negative}
A \emph{negative cycle}\index{negative cycle} is a cycle of negative weight. If one is reachable from $s$ and reaches
$v$, going round it again and again makes $\delta(s, v) = -\infty$.
\end{definition}

\begin{proposition}[Subpaths of shortest paths]\label{prop:b1:shortest-paths:subpath}
Every subpath of a \omterm{def:b1:shortest-paths:shortest}{shortest path} is a \omterm{def:b1:shortest-paths:shortest}{shortest path} between its ends. And if there is no \omterm{def:b1:shortest-paths:negative}{negative cycle}, $\delta(s, v)
\leq \delta(s, u) + w(u, v)$ for every \omterm{def:b1:graph-traversals:graph}{arc} $u \to v$ (the triangle inequality).
\end{proposition}

The proof is \cref{exo:b1:shortest-paths:9}; it is cut and paste again, as in \cref{ch:b1:dynamic-programming}.

\begin{definition}[Relaxation]\label{def:b1:shortest-paths:relaxation}
The algorithms keep an estimate $\code{dist}[v] \geq \delta(s, v)$, with $\code{dist}[s] = 0$ and $+\infty$ elsewhere at
the start. To \emph{relax}\index{relaxation (of an edge)} the \omterm{def:b1:graph-traversals:graph}{arc} $u \to v$ is to set $\code{dist}[v] = \code{dist}[u] +
w(u, v)$, and the parent of $v$ to $u$, if that is smaller.
\end{definition}

Every estimate is the weight of some path, the path of parents, so it never drops below the distance; and once
$\code{dist}[u] = \delta(s, u)$, relaxing $u \to v$ for the last \omterm{def:b1:graph-traversals:graph}{arc} of a \omterm{def:b1:shortest-paths:shortest}{shortest path} to $v$ gives $\code{dist}[v] =
\delta(s, v)$. The algorithms differ only in the order in which they \omterm{def:b1:shortest-paths:relaxation}{relax} \omterm{def:b1:graph-traversals:graph}{arcs}.

\begin{definition}[Shortest-path tree]\label{def:b1:shortest-paths:tree}
A \emph{shortest-path tree}\index{shortest-path tree} from $s$ is a \omterm{def:b1:induction-and-termination:rooted-tree}{rooted tree} on the \omterm{def:b1:graph-traversals:graph}{vertices} reachable from $s$,
rooted at $s$, whose path from $s$ to each \omterm{def:b1:graph-traversals:graph}{vertex} $v$ is a \omterm{def:b1:shortest-paths:shortest}{shortest path}.
\end{definition}

\section{Dijkstra's algorithm with a heap}\label{sec:b1:shortest-paths:dijkstra}

\begin{method}[Dijkstra's algorithm]\label{met:b1:shortest-paths:dijkstra}
With non-negative weights: keep a set of \emph{settled} \omterm{def:b1:graph-traversals:graph}{vertices}, at first empty. Repeatedly settle the unsettled \omterm{def:b1:graph-traversals:graph}{vertex}
of smallest estimate, and \omterm{def:b1:shortest-paths:relaxation}{relax} every \omterm{def:b1:graph-traversals:graph}{arc} leaving it.
\end{method}

A \omterm{def:b1:heaps:binary-heap}{binary heap} gives the smallest estimate. When an estimate decreases, the program pushes a new entry rather than
moving the old one (\emph{lazy deletion}); the stale entry is skipped when it comes off the heap.

\omcode{code/bachelor-1/25-shortest-paths/listings/dijkstra.py}{11}{34}{Dijkstra's algorithm with a \omterm{def:b1:heaps:binary-heap}{binary heap} and
lazy deletion}\label{lst:b1:shortest-paths:dijkstra}

\begin{theorem}[Dijkstra's invariant]\label{thm:b1:shortest-paths:dijkstra}
With non-negative weights, every \omterm{def:b1:graph-traversals:graph}{vertex} $u$ has $\code{dist}[u] = \delta(s, u)$ when it is settled. So at the end
\code{dist} holds the distances, and the parents form a \omterm{def:b1:shortest-paths:tree}{shortest-path tree}.
\end{theorem}

The proof is \cref{exo:b1:shortest-paths:10}: if $u$ were the first \omterm{def:b1:graph-traversals:graph}{vertex} settled with a wrong estimate, a \omterm{def:b1:shortest-paths:shortest}{shortest
path} to $u$ would leave the settled set somewhere, at an \omterm{def:b1:graph-traversals:graph}{arc} $x \to y$, and $y$ would have an estimate at most
$\delta(s, u)$, smaller than $u$'s.

\begin{omfigure}
\begin{tikzpicture}[v/.style={circle, draw, minimum size=6mm, inner sep=0pt, font=\small},
  done/.style={v, fill=omDef!25}, >=stealth, font=\scriptsize]
\foreach \dx/\title/\ds/\da/\db/\dc/\dt/\ss/\sa/\sb in {0/after s/0/4/1/$\infty$/$\infty$/done/v/v,
  50/after b/0/3/1/6/$\infty$/done/v/done, 100/after a/0/3/1/4/9/done/done/done} {
\begin{scope}[xshift=\dx mm]
\node[\ss] (s) at (0,0) {s};
\node[\sa] (a) at (1.6,0.9) {a}; \node[\sb] (b) at (1.6,-0.9) {b}; \node[v] (c) at (2.9,0) {c};
\node[v] (t) at (4.0,0.9) {t};
\draw[->] (s) -- node[above left=-1pt] {4} (a); \draw[->] (s) -- node[below left=-1pt] {1} (b);
\draw[->] (b) -- node[left] {2} (a); \draw[->] (a) -- node[above right=-2pt] {1} (c);
\draw[->] (b) -- node[below right=-2pt] {5} (c); \draw[->] (c) -- node[below right=-1pt] {3} (t);
\draw[->] (a) -- node[above] {6} (t);
\node[omDef] at (0,-0.55) {\ds}; \node[omDef] at (1.6,1.45) {\da}; \node[omDef] at (1.6,-1.45) {\db};
\node[omDef] at (2.9,-0.55) {\dc}; \node[omDef] at (4.0,1.45) {\dt};
\node at (2,-2.0) {\small \title};
\end{scope}}
\end{tikzpicture}
\par\smallskip{\small Dijkstra's algorithm from \code{s}, after settling \code{s}, then \code{b}, then \code{a}; the
estimates are written beside the \omterm{def:b1:graph-traversals:graph}{vertices}. Settling \code{b}, at 1, lowers \code{a} from 4 to 3; settling \code{a}
lowers \code{c} from 6 to 4. Then \code{c} is settled at 4 and \code{t} at 7.}
\end{omfigure}

Each \omterm{def:b1:graph-traversals:graph}{vertex} is settled once and each \omterm{def:b1:graph-traversals:graph}{arc} relaxed once, so at most $m$ entries are pushed, plus the first; each push or
pop costs $O(\log m) = O(\log n)$, since $m \leq n^2$. The total is $O((n + m) \log n)$. A heap with
\emph{\omterm{def:b1:heaps:decrease-key}{decrease-key}} (\cref{ch:b1:heaps}) keeps one entry per \omterm{def:b1:graph-traversals:graph}{vertex}, at the same cost, and less memory.

\omoutput{code/bachelor-1/25-shortest-paths/outputs/dijkstra.txt}{The fastest and the shortest way from the Town Square
to the Hospital; ``pops'' counts the entries taken off the heap, stale ones included.}\label{out:b1:shortest-paths:dijkstra}

\begin{omfigure}
\begin{tikzpicture}
\begin{axis}[omrecord, width=10cm, height=8.6cm, xlabel={$x$ (km)}, ylabel={$y$ (km)}, xmin=-0.2, xmax=6, ymin=-0.2,
    ymax=5.1, axis equal image, grid=none]
  \addplot[black!25, line width=0.6pt] table[x=x, y=y]{code/bachelor-1/25-shortest-paths/figures/roads.dat};
  \addplot[omThm, line width=1.6pt, dashed] table[x=x, y=y]{code/bachelor-1/25-shortest-paths/figures/short.dat};
  \addplot[omDef, line width=1.6pt] table[x=x, y=y]{code/bachelor-1/25-shortest-paths/figures/fast.dat};
  \addplot[only marks, mark=*, mark size=1.3pt, black!60] table[x=x, y=y]
    {code/bachelor-1/25-shortest-paths/figures/places.dat};
  \node[font=\scriptsize, anchor=north west] at (axis cs:4.35,1.40) {Town Square};
  \node[font=\scriptsize, anchor=east] at (axis cs:1.28,4.875) {Hospital};
\end{axis}
\end{tikzpicture}
\par\smallskip{\small The fastest route (solid, 445 s) and the shortest one (dashed, 6270 m), drawn from the data by
\code{figures/routes.py}. They share their first two roads; then the fast one turns west through the Station onto
faster roads, while the short one heads north.}
\end{omfigure}

\section{Negative weights and Bellman--Ford}\label{sec:b1:shortest-paths:bellman-ford}

Weights may be negative: a reward, a gain of height, a refund. Dijkstra's algorithm then fails, because a \omterm{def:b1:graph-traversals:graph}{vertex} can be
settled before a cheaper path reaches it through a negative \omterm{def:b1:graph-traversals:graph}{arc}.

\begin{omfigure}
\begin{tikzpicture}[v/.style={circle, draw, minimum size=6mm, inner sep=0pt, font=\small}, >=stealth,
  font=\scriptsize]
\node[v] (s) at (0,0) {s}; \node[v] (a) at (2,0.8) {a}; \node[v] (b) at (2,-0.8) {b}; \node[v] (t) at (4,0.8) {t};
\draw[->] (s) -- node[above] {2} (a); \draw[->] (s) -- node[below] {3} (b);
\draw[->, omThm, thick] (b) -- node[right] {$-2$} (a); \draw[->] (a) -- node[above] {1} (t);
\end{tikzpicture}
\par\smallskip{\small Dijkstra settles \code{a} at 2 and \omterm{def:b1:shortest-paths:relaxation}{relaxes} $a \to t$, giving \code{t} the estimate 3. Settling
\code{b} at 3, the \omterm{def:b1:graph-traversals:graph}{arc} $b \to a$ lowers \code{a} to 1, but \code{a} is settled and its \omterm{def:b1:graph-traversals:graph}{arcs} are not relaxed again:
\code{t} keeps 3, though its distance is 2.}
\end{omfigure}

\omoutput{code/bachelor-1/25-shortest-paths/outputs/dijkstra_neg.txt}{The listing of \cref{lst:b1:shortest-paths:dijkstra}
on this \omterm{def:b1:graph-traversals:graph}{graph}.}\label{out:b1:shortest-paths:negative}

\begin{method}[Bellman--Ford algorithm]\label{met:b1:shortest-paths:bellman-ford}
\omterm{def:b1:shortest-paths:relaxation}{Relax} every \omterm{def:b1:graph-traversals:graph}{arc}, in any order; repeat, in \emph{rounds}, until a round changes nothing. If round $n$ still changes an
estimate, report a \omterm{def:b1:shortest-paths:negative}{negative cycle}.
\end{method}

\omcode{code/bachelor-1/25-shortest-paths/listings/bellman_ford.py}{9}{22}{Bellman--Ford, with a stop at the first
quiet round}\label{lst:b1:shortest-paths:bellman-ford}

\begin{proposition}[Rounds of Bellman--Ford]\label{prop:b1:shortest-paths:rounds}
If no \omterm{def:b1:shortest-paths:negative}{negative cycle} is reachable from $s$, then after $k$ rounds $\code{dist}[v] = \delta(s, v)$ for every \omterm{def:b1:graph-traversals:graph}{vertex} $v$
that has a \omterm{def:b1:shortest-paths:shortest}{shortest path} of at most $k$ \omterm{def:b1:graph-traversals:graph}{arcs}; so $n - 1$ rounds suffice and round $n$ changes nothing. If a \omterm{def:b1:shortest-paths:negative}{negative
cycle} is reachable from $s$, every round changes some estimate.
\end{proposition}

\begin{proof}
Without a reachable \omterm{def:b1:shortest-paths:negative}{negative cycle}, every reachable $v$ has a \omterm{def:b1:shortest-paths:shortest}{shortest path} that is simple, so of at most $n - 1$ \omterm{def:b1:graph-traversals:graph}{arcs}.
By induction on $k$: if $s = v_0, \dots, v_k = v$ is a \omterm{def:b1:shortest-paths:shortest}{shortest path}, $\code{dist}[v_{k-1}] = \delta(s, v_{k-1})$ after
round $k - 1$ (its prefix is a \omterm{def:b1:shortest-paths:shortest}{shortest path}, \cref{prop:b1:shortest-paths:subpath}), and round $k$ \omterm{def:b1:shortest-paths:relaxation}{relaxes} $v_{k-1}
\to v$. Estimates never go below distances, so they stop changing. Conversely, if a round changes nothing, then
$\code{dist}[v] \leq \code{dist}[u] + w(u, v)$ for every \omterm{def:b1:graph-traversals:graph}{arc}; summing along a reachable cycle, where every estimate is
finite, gives $0 \leq$ the weight of the cycle: no \omterm{def:b1:shortest-paths:negative}{negative cycle} is reachable.
\end{proof}

Each round costs $O(n + m)$, so the whole costs $O(nm)$, far more than Dijkstra's algorithm, but it accepts any weights.

\omoutput{code/bachelor-1/25-shortest-paths/outputs/bellman_ford.txt}{Five rounds suffice on the town, the last one
quiet; the \omterm{def:b1:graph-traversals:graph}{arc} $t \to b$ closes the cycle $b \to a \to t \to b$ of weight $-1$.}\label{out:b1:shortest-paths:bellman-ford}

Bellman published the method in 1958 as an algorithm for a routing problem; it also carries the name of Lester Ford,
who studied the same kind of problem.

\section{Shortest paths in a DAG}\label{sec:b1:shortest-paths:dag}

In a \omterm{def:b1:graph-traversals:dag}{directed acyclic graph}, relaxing the \omterm{def:b1:graph-traversals:graph}{arcs} leaving each \omterm{def:b1:graph-traversals:graph}{vertex}, the \omterm{def:b1:graph-traversals:graph}{vertices} taken in a \omterm{def:b1:graph-traversals:topo}{topological order}, gives
all distances from $s$ in one pass, whatever the signs of the weights: when $u$ is reached, every \omterm{def:b1:graph-traversals:graph}{arc} entering $u$ has
been relaxed, and by induction along the order $\code{dist}[u] = \delta(s, u)$. The cost is $O(n + m)$, the
\omterm{def:b1:graph-traversals:topo}{topological order} included. The same pass with maxima in place of minima gives \emph{longest} paths, a problem that is
hard in general \omterm{def:b1:graph-traversals:graph}{graphs} but easy without cycles.

\begin{definition}[Critical path]\label{def:b1:shortest-paths:critical}
A project is a set of tasks with durations and prerequisites, which form a DAG. The earliest start of a task is the
largest total duration of a chain of tasks that must precede it; a \emph{critical path}\index{critical path} is a chain
of tasks of largest total duration: its length is the shortest possible duration of the project.
\end{definition}

\omcode{code/bachelor-1/25-shortest-paths/listings/dag_paths.py}{32}{45}{Earliest starts and a \omterm{def:b1:shortest-paths:critical}{critical path}, in
\omterm{def:b1:graph-traversals:topo}{topological order}}\label{lst:b1:shortest-paths:dag}

\omoutput{code/bachelor-1/25-shortest-paths/outputs/dag_paths.txt}{A small software project: delay any task of the
\omterm{def:b1:shortest-paths:critical}{critical path} and the launch is delayed.}\label{out:b1:shortest-paths:dag}

The manual can start on day 7 and has until day 9, two days before training must start on day 11: it has two days of
slack. The tasks of
the \omterm{def:b1:shortest-paths:critical}{critical path} have none.

\section{All pairs: Floyd--Warshall}\label{sec:b1:shortest-paths:floyd}

For the distances between all pairs, one can run Dijkstra's algorithm from every \omterm{def:b1:graph-traversals:graph}{vertex}, at a cost of $O(n (n + m) \log
n)$. A dynamic programme on the adjacency matrix is shorter to write and handles negative weights.

\begin{method}[Floyd--Warshall algorithm]\label{met:b1:shortest-paths:floyd}
Number the \omterm{def:b1:graph-traversals:graph}{vertices} $0, \dots, n - 1$. Let $D_k(i, j)$ be the smallest weight of a path from $i$ to $j$ whose
intermediate \omterm{def:b1:graph-traversals:graph}{vertices} are all less than $k$. Then $D_0(i, j)$ is $0$ if $i = j$, $w(i, j)$ for an \omterm{def:b1:graph-traversals:graph}{arc} and $+\infty$
otherwise, and
\[ D_{k+1}(i, j) = \min\bigl(D_k(i, j),\; D_k(i, k) + D_k(k, j)\bigr). \]
The table is updated \omterm{def:b1:sorting:comparison-sort}{in place}, $k$ in the outer loop; $D_n$ is the distance table.
\end{method}

\omcode{code/bachelor-1/25-shortest-paths/listings/floyd.py}{9}{36}{Floyd--Warshall, with a table of next \omterm{def:b1:graph-traversals:graph}{vertices} to
rebuild the paths}\label{lst:b1:shortest-paths:floyd}

The triple loop costs $\Theta(n^3)$ and the tables $\Theta(n^2)$. On a dense \omterm{def:b1:graph-traversals:graph}{graph}, where $m$ is close to $n^2$, it beats
$n$ runs of Dijkstra by a logarithm and more; on a sparse one, such as roads, $n$ runs of Dijkstra cost about $n^2
\log n$ and win when $n$ is large. Updating \omterm{def:b1:sorting:comparison-sort}{in place} is safe: in round $k$, the entries $D(i, k)$ and $D(k, j)$ do not
change, since a path to or from $k$ gains nothing by going through $k$ again, without \omterm{def:b1:shortest-paths:negative}{negative cycles}.

\omoutput{code/bachelor-1/25-shortest-paths/outputs/floyd.txt}{All pairs in the town; the one-way road makes the
return from the Library longer.}\label{out:b1:shortest-paths:floyd}

Floyd published the algorithm in 1962 as a short procedure for \omterm{def:b1:shortest-paths:shortest}{shortest paths}; the same year, Warshall proved the
same loop correct for the reachability relation, with ``or'' and ``and'' in place of ``min'' and ``$+$''.

\begin{definition}[Eccentricity, centre, diameter]\label{def:b1:shortest-paths:centre}
In a \omterm{def:b1:graph-traversals:connected}{strongly connected} \omterm{def:b1:graph-traversals:weighted}{weighted graph}, the \emph{eccentricity}\index{eccentricity} of a \omterm{def:b1:graph-traversals:graph}{vertex} is its greatest
distance to another \omterm{def:b1:graph-traversals:graph}{vertex}. A \emph{centre}\index{centre (of a graph)} is a \omterm{def:b1:graph-traversals:graph}{vertex} of smallest eccentricity, the
\emph{radius}; the \emph{diameter}\index{diameter (of a graph)} is the greatest eccentricity, the largest distance
between two \omterm{def:b1:graph-traversals:graph}{vertices}.
\end{definition}

\section{Exercises}\label{sec:b1:shortest-paths:exercises}

\begin{exercise}[$\star$]\label{exo:b1:shortest-paths:1}
In the \omterm{def:b1:graph-traversals:graph}{graph} drawn for Dijkstra's algorithm, give the order in which the \omterm{def:b1:graph-traversals:graph}{vertices} are settled and the final distances, starting
from \code{a} instead of \code{s}.
\end{exercise}

\begin{exercise}[$\star$]\label{exo:b1:shortest-paths:2}
Run Bellman--Ford from $s$ on the \omterm{def:b1:graph-traversals:graph}{arcs} $a \to c$ (2), $b \to a$ ($-1$), $s \to a$ (4), $s \to b$ (2), $b \to c$ (5),
relaxed in this order in every round. Give the estimates after each round, and the number of rounds.
\end{exercise}

\begin{exercise}[$\star$]\label{exo:b1:shortest-paths:3}
Three \omterm{def:b1:graph-traversals:graph}{vertices} 0, 1, 2 have \omterm{def:b1:graph-traversals:graph}{arcs} $0 \to 1$ (5), $1 \to 2$ (2), $0 \to 2$ (9), $2 \to 0$ (1). Give the matrices $D_0$, $D_1$
and $D_2$ of Floyd--Warshall.
\end{exercise}

\begin{exercise}[$\star$]\label{exo:b1:shortest-paths:4}
In the project of \cref{out:b1:shortest-paths:dag}, the database now takes 3 days. What are the new length and \omterm{def:b1:shortest-paths:critical}{critical
path}?
\end{exercise}

\begin{exercise}[$\star\star$]\label{exo:b1:shortest-paths:5}
This version fixes a \omterm{def:b1:graph-traversals:graph}{vertex}'s distance when it is first pushed. Explain the output and the difference with
\cref{lst:b1:shortest-paths:dijkstra}.
\omcode{code/bachelor-1/25-shortest-paths/listings/exo_settle_push.py}{6}{14}{Fixed when pushed}
\omoutput[2-2]{code/bachelor-1/25-shortest-paths/outputs/exo_settle_push.txt}{\omterm{def:b1:graph-traversals:graph}{Arcs} $s \to a$ (5), $s \to b$ (1), $b \to
a$ (1).}
\end{exercise}

\begin{exercise}[$\star\star$]\label{exo:b1:shortest-paths:6}
With the loop on $k$ innermost, 350 entries of the town's table are wrong. Why?
\omcode{code/bachelor-1/25-shortest-paths/listings/exo_floyd_inner.py}{17}{21}{The loops in another order}
\omoutput[2-2]{code/bachelor-1/25-shortest-paths/outputs/exo_floyd_inner.txt}{Against Floyd--Warshall.}
\end{exercise}

\begin{exercise}[$\star\star$]\label{exo:b1:shortest-paths:7}
Compute the smallest weight of a path of at most $k$ \omterm{def:b1:graph-traversals:graph}{arcs} from $s$ to every \omterm{def:b1:graph-traversals:graph}{vertex}, with any weights, in $O(k(n + m))$,
and test it.
\end{exercise}

\begin{exercise}[$\star\star$]\label{exo:b1:shortest-paths:8}
Each \omterm{def:b1:graph-traversals:graph}{arc} has a capacity, for instance the height limit of a road. Find, from $s$, the paths whose smallest capacity is
as large as possible, by a variant of Dijkstra's algorithm, and justify it.
\end{exercise}

\begin{exercise}[$\star\star$]\label{exo:b1:shortest-paths:9}
Prove \cref{prop:b1:shortest-paths:subpath}.
\end{exercise}

\begin{exercise}[$\star\star\star$]\label{exo:b1:shortest-paths:10}
Prove \cref{thm:b1:shortest-paths:dijkstra}.
\end{exercise}

\begin{exercise}[$\star\star\star$]\label{exo:b1:shortest-paths:11}
In a weighted DAG, count the \omterm{def:b1:shortest-paths:shortest}{shortest paths} from $s$ to every \omterm{def:b1:graph-traversals:graph}{vertex} in $O(n + m)$, and prove it.
\end{exercise}

\begin{exercise}[$\star\star\star$]\label{exo:b1:shortest-paths:12}
Write a function that returns a \omterm{def:b1:shortest-paths:negative}{negative cycle} of a \omterm{def:b1:graph-traversals:graph}{graph}, as a list of \omterm{def:b1:graph-traversals:graph}{vertices}, or \code{None}, in $O(nm)$.
\end{exercise}

\section{Problem: The centre of Brenmoor}\label{sec:b1:shortest-paths:problem}

\begin{problem}[{Weekend problem --- The centre of Brenmoor}]\label{pb:b1:shortest-paths:1}
Where in Brenmoor should the new ambulance station go, so that the farthest place is as close as possible? Travel times
are those of \cref{lst:b1:shortest-paths:wgraph}.
\omcode{code/bachelor-1/25-shortest-paths/listings/centre.py}{8}{9}{Eccentricities from the distance
table}\label{lst:b1:shortest-paths:centre}
\omoutput{code/bachelor-1/25-shortest-paths/outputs/centre.txt}{The three most central places, the two least central,
the \omterm{def:b1:shortest-paths:centre}{centre} and the \omterm{def:b1:shortest-paths:centre}{diameter}.}\label{out:b1:shortest-paths:centre}

\needspace{4\baselineskip}
\medskip\noindent\textbf{Part I --- Modelling.}
\begin{enumerate}
\item Compute the travel time of a road of 1360 m at 50 km/h, as the program does.
\item Why does the Station have an \omterm{def:b1:graph-traversals:graph}{arc} of 120 s to the Town Hall but only 35 s to River Park?
\item How many \omterm{def:b1:graph-traversals:graph}{arcs} does the town have, and why not 106?
\item Why are the weights in whole seconds, and what does rounding cost on a route of seven roads?
\item Why must a one-way road be an \omterm{def:b1:graph-traversals:graph}{arc}, not an edge, in this model?
\item Could the weights be negative here? Which algorithms would still apply?
\end{enumerate}

\needspace{4\baselineskip}
\medskip\noindent\textbf{Part II --- Dijkstra from the Town Square.}
\begin{enumerate}[resume]
\item Read the fastest and the shortest routes to the Hospital, and compare their lengths and times.
\item Why did the search pop 35 entries for 30 places?
\item Bound the number of pops in terms of $n$ and $m$.
\item Why do the parents form a tree, and why is it a \omterm{def:b1:shortest-paths:tree}{shortest-path tree}?
\item Up to which place do the two routes agree, and what does \cref{prop:b1:shortest-paths:subpath} say of their first parts?
\item \omterm{def:b1:graph-traversals:bfs}{Breadth-first search} found a route of six roads to the Hospital in \cref{out:b1:graph-traversals:bfs}: compare it with the two routes.
\end{enumerate}

\needspace{4\baselineskip}
\medskip\noindent\textbf{Part III --- All pairs.}
\begin{enumerate}[resume]
\item Prove by induction on $k$ that $D_k(i, j)$, as defined by the recurrence, is the smallest weight of a path from
$i$ to $j$ with all intermediate \omterm{def:b1:graph-traversals:graph}{vertices} less than $k$.
\item Deduce that $D_n$ is the distance table.
\item Show that the in-place update computes the same table.
\item Explain how \code{nxt} rebuilds a route, and why \code{nxt[i][j] = nxt[i][k]} is right.
\item Why do the Library and the Town Square have different distances in the two directions?
\item Compare the costs of Floyd--Warshall and of 30 runs of Dijkstra on the town.
\end{enumerate}

\needspace{4\baselineskip}
\medskip\noindent\textbf{Part IV --- The centre.}
\begin{enumerate}[resume]
\item Explain \code{max(d[i])}: why is the distance from a place to itself harmless?
\item Read the \omterm{def:b1:shortest-paths:centre}{centre}, the radius and the \omterm{def:b1:shortest-paths:centre}{diameter}.
\item Why does the Post Office appear as farthest for several places?
\item Show that the \omterm{def:b1:shortest-paths:centre}{diameter} is at most twice the radius when every road is two-way.
\item Does that bound hold for Brenmoor, with its one-way road?
\item State the named result, \emph{the town's \omterm{def:b1:shortest-paths:centre}{centre} and \omterm{def:b1:shortest-paths:centre}{diameter}}, with its proof.
\end{enumerate}
\end{problem}
